\documentclass{report}
\usepackage{amsmath,amssymb}
\setlength{\parindent}{0mm}
\setlength{\parskip}{1em}
\begin{document}
\begin{center}
\rule{6in}{1pt} \
{\large Ben Bergen \\
{\bf Task-Graph and Functional Programming Models: The New Paradigm}}

P O Box 1663 \\ MS B287 \\ Los Alamos \\ NM 87545
\\
{\tt bergen@lanl.gov}\end{center}

The Message Passing Interface (MPI) is an example of a distributed-memory
communication model that has served us well through the CISC processor
era. However, because of MPI�s low-level interface, which requires the
user to manage raw memory buffers, and its bulk-synchronous communication
model, MPI will have great difficulty in scaling to exascale systems and
beyond. Additionally, the MPI model cannot be easily extended to include
the fault tolerance and resilience features that will be required to run
at scale on modern computing architectures. A new approach is needed.

A trend that is gaining momentum in the computer and computational
science communities is the use of data-driven models that employ
task-graph runtimes to map data and tasks onto a distributed system. This
approach offers a higher level of abstraction that frees the user from
explicitly knowing the layout and location of the data regions accessed
by a task. Task and data dependencies can be represented as a directed
acyclic graph (DAG), allowing heuristics and auto-tuning strategies to be
applied at runtime to optimize the execution of the DAG. This approach
shows great promise. However, it will require a new understanding of how
our existing numerical methods can be mapped into this model. At the
same, there is an opportunity for us to influence the features and
requirements for implementations of task-graph models.

Although there are several efforts underway to develop useful task-graph
implementations, the Legion project at Stanford University is a clear
front-runner. This presentation is part one of a two-part series designed
to introduce some of the Legion concepts to the audience. Here we discuss
high-level task-graph concepts and provide an example of how a geometric
multigrid algorithm might be implemented using a task-graph runtime.


\end{document}
